\documentclass[12pt,a4paper]{article}
\usepackage[utf8]{vietnam}
\usepackage{amsmath,amssymb}
\usepackage{geometry}
\geometry{top=1.5cm,bottom=1.2cm,left=1.4cm,right=1.4cm,headheight=28pt,headsep=12pt,footskip=18pt}
\usepackage{tikz}
\usepackage{xcolor}
\usepackage{enumerate}
\usepackage{graphicx}
\usepackage[version=4]{mhchem}
\usepackage{fancyhdr}
\usepackage{ex_test}

\renewcommand{\baselinestretch}{0.95}
\setlength{\parskip}{0pt}

\definecolor{hfRed}{RGB}{211,47,47}
\definecolor{hfGreen}{RGB}{139,195,144}

\pagestyle{fancy}
\fancyhf{}
\fancyhead[L]{\small\itshape\color{hfRed} Tài liệu lưu hành nội bộ}
\fancyhead[R]{\small\itshape\color{hfRed} Thầy \textbf{TRẦN VĂN THẠNH} - 0777.470.803}
\renewcommand{\headrulewidth}{0.8pt}
\renewcommand{\headrule}{\hbox to\headwidth{\color{hfGreen}\leaders\hrule height \headrulewidth\hfill}}

\fancyfoot[L]{\small\itshape\color{hfRed} CS1: 6/15 Nguyễn Hoàng, P. Kim Long, TP Huế \quad|\quad CS2: 24 Đặng Thái Thân, TP Huế}
\fancyfoot[R]{\small\bfseries\color{hfRed} Trang \thepage}
\renewcommand{\footrulewidth}{0.8pt}
\renewcommand{\footrule}{\hbox to\headwidth{\color{hfGreen}\leaders\hrule height \footrulewidth\hfill}}

\fancypagestyle{firstpage}{
  \fancyhf{}
  \renewcommand{\headrulewidth}{0pt}
  \fancyfoot[L]{\small\itshape\color{hfRed} CS1: 6/15 Nguyễn Hoàng, P. Kim Long, TP Huế \quad|\quad CS2: 24 Đặng Thái Thân, TP Huế}
  \fancyfoot[R]{\small\bfseries\color{hfRed} Trang \thepage}
  \renewcommand{\footrulewidth}{0.8pt}
  \renewcommand{\footrule}{\hbox to\headwidth{\color{hfGreen}\leaders\hrule height \footrulewidth\hfill}}
}

\begin{document}
\thispagestyle{firstpage}

\noindent
\begin{minipage}[t]{0.42\textwidth}
	\centering\fontsize{10pt}{12pt}\selectfont
	\textbf{THPT NGUYỄN ĐÌNH CHIỂU}\\
	\textbf{ĐỀ CHÍNH THỨC}\\
	\fbox{\textbf{MÃ ĐỀ: 1011}}
\end{minipage}\hfill
\begin{minipage}[t]{0.56\textwidth}
	\centering\small
	\textbf{ĐỀ KIỂM TRA GIỮA KÌ I NĂM HỌC 2025 - 2026}\\
	\textbf{Môn thi: HÓA HỌC 12}\\
	\textit{Thời gian làm bài: 45 phút (không tính thời gian phát đề)}
\end{minipage}

\smallskip
\noindent\hrulefill
\smallskip

\noindent\textbf{Họ và tên học sinh:} \dotfill\ \textbf{Lớp: 12A1}

\medskip
\noindent
\textbf{A. PHẦN TRẮC NGHIỆM (7 điểm)}

\smallskip
\noindent
\textbf{PHẦN I (3 điểm -- 12 câu): Câu trắc nghiệm nhiều phương án lựa chọn.} \textit{Thí sinh trả lời từ câu 1 đến câu 12. Mỗi câu hỏi thí sinh chỉ chọn 1 phương án.}

\Opensolutionfile{ans}[ans-hs]

\begin{ex}
Cellulose trinitrate được điều chế từ cellulose và nitric acid đặc có xúc tác sulfuric acid đặc, nóng. Để có $44{,}55\text{ kg}$ cellulose trinitrate, cần dùng dung dịch chứa $m\text{ kg}$ nitric acid (hiệu suất phản ứng đạt 90\%). Giá trị của $m$ là
\choice
{$28{,}350\text{ kg}$}
{$21{,}234\text{ kg}$}
{$25{,}515\text{ kg}$}
{\True $31{,}500\text{ kg}$}
\end{ex}

\begin{ex}
Methyl salicylate (chất X) là sản phẩm tự nhiên của rất nhiều loại cây, thường được kết hợp với các loại tinh dầu khác dùng làm thuốc bôi ngoài da, thuốc xoa bóp, cao dán giảm đau, chống viêm. X có công thức cấu tạo như sau:
\begin{center}
	\includegraphics[width=3.8cm]{San_Pham/images_crop/fig_cau2_methyl_salicylate.png}
\end{center}
Cho các phát biểu sau:
\begin{enumerate}[(a)]
	\item Công thức phân tử của X là $\ce{C8H8O3}$.
	\item Phân tử X chứa 31,58\% oxygen về khối lượng.
	\item $a\text{ mol X}$ phản ứng tối đa với $a\text{ mol Na}$, sinh ra $a\text{ mol H2}$.
	\item $a\text{ mol X}$ phản ứng tối đa với $2a\text{ mol NaOH}$.
	\item X là hợp chất hữu cơ tạp chức, chứa đồng thời chức ester và chức alcohol.
\end{enumerate}
Số phát biểu \textbf{sai} là
\choice
{1}
{\True 2}
{4}
{3}
\end{ex}

\begin{ex}
Cho các chất sau: glucose, fructose, maltose, saccharose, cellulose và tinh bột. Số chất tạo phức màu xanh lam với $\ce{Cu(OH)2}$ trong môi trường kiềm là
\choice
{3}
{2}
{\True 4}
{1}
\end{ex}

\begin{ex}
Cho các chất: ethanol, acetic acid, methyl fomate, propionic acid. Chất nào có nhiệt độ sôi thấp nhất?
\choice
{\True methyl fomate}
{acetic acid}
{ethanol}
{propionic acid}
\end{ex}

\begin{ex}
Cho các chất: $\ce{CH3[CH2]14COONa}$, $\ce{CH3[CH2]10CH2OSO3Na}$, $\ce{CH3[CH2]7CH=CH[CH2]7COONa}$, $\ce{CH3[CH2]16COOK}$, $\ce{CH3COONa}$, $\ce{(C15H31COO)3C3H5}$. Số chất có thể là thành phần chính của xà phòng là
\choice
{\True 3}
{4}
{1}
{2}
\end{ex}

\begin{ex}
Để rửa sạch chai lọ đựng dung dịch aniline, nên dùng cách nào sau đây?
\choice
{Rửa bằng dung dịch $\ce{NaOH}$ sau đó rửa lại bằng nước}
{\True Rửa bằng dung dịch $\ce{HCl}$ sau đó rửa lại bằng nước}
{Rửa bằng xà phòng}
{Rửa bằng nước}
\end{ex}

\begin{ex}
Cho sơ đồ chuyển hóa sau:
\[\ce{X_{(l)} + 2NaOH_{(aq)} -> CH2(COONa)2_{(aq)} + CH3OH_{(aq)} + C2H5OH_{(aq)}}\]
Nhận xét nào sau đây \textbf{không đúng} về chất X trong phản ứng trên?
\choice
{Công thức cấu tạo của (X) được biểu diễn như sau:\\
\centerline{\includegraphics[width=6.5cm]{San_Pham/images_crop/fig_cau7_malonate.png}}}
{X là ester no có hai nhóm chức có công thức phân tử $\ce{C6H10O4}$}
{Tên của X là ethyl methyl malonate}
{\True X có nhiệt độ sôi cao vì có liên kết hydrogen giữa các phân tử}
\end{ex}

\begin{ex}
Tiến hành thí nghiệm điều chế ethyl acetate theo các bước sau:
\begin{itemize}
	\item Bước 1: Cho $1\text{ mL } \ce{C2H5OH}$, $1\text{ mL } \ce{CH3COOH}$ và vài giọt dung dịch $\ce{H2SO4}$ đặc vào ống nghiệm.
	\item Bước 2: Lắc đều ống nghiệm, đun cách thủy (trong nồi nước nóng) khoảng 5 - 6 phút ở $65 - 70^\circ\text{C}$.
	\item Bước 3: Làm lạnh, sau đó rót $2\text{ mL}$ dung dịch $\ce{NaCl}$ bão hòa vào ống nghiệm.
\end{itemize}
Phát biểu nào sau đây là \textbf{sai}?
\choice
{$\ce{H2SO4}$ đặc có vai trò vừa làm chất xúc tác vừa làm tăng hiệu suất tạo sản phẩm}
{Sau bước 3, chất lỏng trong ống nghiệm tách thành hai lớp}
{\True Mục đích chính của việc thêm dung dịch $\ce{NaCl}$ bão hòa là để tránh phân hủy sản phẩm}
{Sau bước 2, trong ống nghiệm vẫn còn $\ce{C2H5OH}$ và $\ce{CH3COOH}$}
\end{ex}

\begin{ex}
Cho sơ đồ chuyển hóa:
\[\ce{X (C10H16O7N2) ->[NaOH\ d\text{ư}] Y ->[HCl\ d\text{ư}] Z}\]
Biết X là dipeptide của một $\alpha$-amino acid \textbf{T} có cấu tạo không phân nhánh; mỗi mũi tên ứng với một phương trình hóa học của phản ứng giữa hai chất tương ứng. Phát biểu nào sau đây \textbf{đúng}?
\choice
{Phần trăm khối lượng của nguyên tố chlorine trong phân tử chất \textbf{Z} chiếm 19,452\%}
{\True Ở điều kiện thường, chất \textbf{T} dễ tan trong nước và có nhiệt độ nóng chảy cao}
{Chất \textbf{Y} dùng làm gia vị thức ăn (gọi là mì chính hay bột ngọt)}
{X tác dụng tối đa với dung dịch $\ce{NaOH}$ theo tỉ lệ 1 : 3}
\end{ex}

\begin{ex}
Carbohydrate (X) có công thức cấu tạo dưới đây:
\begin{center}
	\includegraphics[width=8.0cm]{San_Pham/images_crop/fig_cau10_maltose.png}
\end{center}
Nhận định nào sau đây là \textbf{đúng} khi nói về (X)?
\choice
{(X) có thể là saccharose}
{\True (X) còn được gọi là đường mạch nha được sản xuất từ ngũ cốc}
{(X) không có tính khử}
{(X) được cấu tạo từ 1 đơn vị $\alpha$-glucose và 1 đơn vị $\beta$-fructose qua liên kết $\alpha$-1,4-glycoside}
\end{ex}

\begin{ex}
Glutamic acid có vai trò quan trọng trong quá trình xây dựng cấu trúc tế bào của con người. Ngoài ra, muối monosodium glutamate còn được dùng chế biến gia vị thức ăn (bột ngọt hay mì chính). Glutamic acid có cấu trúc như hình vẽ bên dưới và có điểm đẳng điện $\text{pI} = 3{,}2$ ($\text{pI}$ là giá trị pH mà khi đó amino acid có nồng độ ion lưỡng cực là cực đại. Khi $\text{pH} < \text{pI}$ thì amino acid đó tồn tại chủ yếu ở dạng cation, còn khi $\text{pH} > \text{pI}$ thì amino acid đó tồn tại chủ yếu ở dạng anion).
\begin{center}
	\includegraphics[width=6.0cm]{San_Pham/images_crop/fig_cau11_glutamic_acid.png}
\end{center}
Cho các phát biểu sau:
\begin{enumerate}[(a)]
	\item Glutamic acid thuộc loại hợp chất hữu cơ tạp chức, trong phân tử chứa hai loại nhóm chức.
	\item Tên thay thế của glutamic acid là 2-aminopentane-1,5-dioic acid.
	\item Trong dung dịch $\text{pH} = 3{,}2$, glutamic acid tồn tại chủ yếu ở dạng $\ce{HOOC-CH2-CH2-CH(NH2)-COO^-}$.
	\item Trong dung dịch $\text{pH} = 6$, có thể tách hỗn hợp gồm glutamic acid và lysine ($\text{pI} = 9{,}7$) bằng phương pháp điện di.
\end{enumerate}
Số phát biểu \textbf{đúng} là
\choice
{1}
{4}
{2}
{\True 3}
\end{ex}

\begin{ex}
Cho các phát biểu sau:
\begin{enumerate}[(a)]
	\item Chất béo được gọi chung là triglyceride.
	\item Chất béo nhẹ hơn nước, không tan trong nước nhưng tan nhiều trong dung môi hữu cơ.
	\item Phản ứng thủy phân chất béo trong môi trường acid là phản ứng thuận nghịch.
	\item Ở nhiệt độ thường, chất béo chứa nhiều gốc acid béo không no thường ở thể rắn.
\end{enumerate}
Số phát biểu \textbf{đúng} là
\choice
{1}
{\True 3}
{2}
{4}
\end{ex}

\Closesolutionfile{ans}

\vspace{0.3cm}
\noindent
\textbf{PHẦN II (3 điểm -- 3 câu): Câu trắc nghiệm đúng sai.} \textit{Thí sinh trả lời từ câu 1 đến câu 3. Trong mỗi ý a), b), c), d) ở mỗi câu thí sinh chọn đúng hoặc sai.}

\setcounter{ex}{0}

\begin{ex}
Triglyceride đóng vai trò là nguồn cung cấp năng lượng và chuyên chở các chất béo trong quá trình trao đổi chất. Cho triglyceride X có công thức cấu tạo như hình sau:
\begin{center}
	\includegraphics[width=13.5cm]{San_Pham/images_crop/fig_cau1_p2_triglyceride.png}
\end{center}
\begin{enumerate}[a)]
	\item X được cấu tạo từ các acid béo là palmitic acid, oleic acid và linoleic acid.
	\item Thủy phân hoàn toàn $427\text{ kg}$ triglyceride X bằng lượng $\ce{NaOH}$ dư thu được $441\text{ kg}$ muối của acid béo.
	\item Trong các acid béo cấu tạo nên X có một acid béo thuộc loại acid béo omega -8.
	\item Tổng số nguyên tử trong 1 phân tử X là 158.
\end{enumerate}
\loigiai{}
\end{ex}

\begin{ex}
Các peptide có phản ứng thủy phân trong môi trường acid và môi trường kiềm, ngoài ra các peptide có từ 2 liên kết peptide trở lên phản ứng với $\ce{Cu(OH)2}$ trong môi trường kiềm tạo thành phức chất màu tím đặc trưng, gọi là phản ứng màu biuret.
\begin{enumerate}[a)]
	\item Dung dịch của các polypeptide hoà tan $\ce{Cu(OH)2}$ cho dung dịch có màu xanh tím.
	\item Thủy phân hoàn toàn $0{,}1\text{ mol}$ Glu--Ala--Lys cần vừa đủ $300\text{ mL}$ dung dịch $\text{KOH 1M}$.
	\item Thủy phân hoàn toàn $2{,}17\text{ gam}$ tripeptide mạch hở X (được tạo nên từ hai $\alpha$-amino acid có công thức dạng $\ce{H2NC_xH_yCOOH}$) bằng dung dịch $\ce{NaOH}$ dư, thu được $3{,}19\text{ gam}$ muối. Mặt khác, thủy phân hoàn toàn $3{,}255\text{ gam}$ X bằng dung dịch $\ce{HCl}$ dư, thu được $5{,}4375\text{ gam}$ muối.
	\item Gly--Ala không có phản ứng màu biuret với $\ce{Cu(OH)2}$.
\end{enumerate}
\loigiai{}
\end{ex}

\begin{ex}
Khi nấu rượu hoặc làm bánh mì từ bột bánh mì có chứa men (một loại nấm) và đường. Khi bột được để ở nơi ấm áp, các tế bào nấm men sẽ ăn đường để lấy năng lượng. Enzyme trong nấm men xúc tác cho phản ứng được gọi là quá trình lên men sinh ra ethanol và khí carbon dioxide:
\[\ce{(C6H10O5)_n} \xrightarrow[\text{enzyme}]{+\ce{H2O}} \ce{n C6H12O6} \xrightarrow{\text{enzyme}} \ce{2n C2H5OH + 2n CO2}\]
\begin{enumerate}[a)]
	\item Để điều chế $5\text{ lít}$ ethyl alcohol $46^\circ$ cần $5{,}4\text{ kg}$ bột bánh mì trên (chứa 75\% tinh bột, còn lại là tạp chất trơ). Biết hiệu suất của cả quá trình là 80\% và khối lượng riêng của ethyl alcohol nguyên chất là $0{,}8\text{ g/mL}$.
	\item Nếu đun nóng hỗn hợp lên men để tách $\ce{C2H5OH}$, đó là phương pháp chưng cất dựa trên sự khác nhau về độ tan.
	\item Lên men rượu luôn xảy ra tối ưu ở khoảng $30 - 35^\circ\text{C}$, nếu nhiệt độ thấp nấm men sẽ chết.
	\item Nếu cung cấp nhiều oxy cho môi trường lên men, hiệu suất tạo ethanol sẽ tăng.
\end{enumerate}
\loigiai{}
\end{ex}

\vspace{0.3cm}
\noindent
\textbf{PHẦN III (1 điểm -- 4 câu): Câu trắc nghiệm yêu cầu trả lời ngắn.} \textit{Thí sinh trả lời từ câu 1 đến câu 4.}

\setcounter{ex}{0}

\begin{ex}
Một loại chất béo có chứa 75\% tristearine về khối lượng. Xà phòng hóa hoàn toàn $17{,}8\text{ kg}$ chất béo này trong dung dịch $\text{KOH}$, đun nóng thu được $x$ bánh xà phòng. Biết rằng trong mỗi bánh xà phòng có chứa $63\text{ gam}$ potassium stearate. Giá trị của $x$ là bao nhiêu?
\loigiai{}
\end{ex}

\begin{ex}
Aspirin là một hợp chất được sử dụng làm giảm đau, hạ sốt được điều chế theo phản ứng sau:
\[\ce{(CH3CO)2O + HOC6H4COOH <=> CH3COOC6H4COOH + CH3COOH}\]
Để sản xuất 500 viên thuốc aspirin cần tối thiểu $20{,}7\text{ gam}$ salicylic acid. Biết rằng mỗi viên thuốc có chứa $y\text{ mg}$ aspirin và hiệu suất phản ứng đạt 75\%. Giá trị của $y$ là bao nhiêu?
\loigiai{}
\end{ex}

\begin{ex}
Cho các nhận định sau:
\begin{enumerate}[(a)]
	\item Protein dạng hình cầu không tan trong nước, còn dạng hình sợi tan trong nước.
	\item Một trong những tính chất hóa học đặc trưng của protein là phản ứng thủy phân.
	\item Phản ứng của protein với nitrous acid cho sản phẩm màu vàng.
	\item Sự đông tụ không làm thay đổi cấu tạo ban đầu của protein bị biến đổi.
	\item Trong cơ thể, enzyme không đóng vai trò là chất xúc tác sinh học.
\end{enumerate}
Số phát biểu đúng là
\loigiai{}
\end{ex}

\begin{ex}
Để tráng một số lượng gương soi có diện tích bề mặt $35\text{ dm}^2$ với độ dày $0{,}08\text{ }\mu\text{m}$ người ta đun nóng dung dịch chứa $30{,}6\text{ gam}$ glucose với một lượng dung dịch $\ce{AgNO3}$ trong amoniac. Biết khối lượng riêng của silver là $10{,}49\text{ kg/L}$, hiệu suất phản ứng tráng gương là 80\% (tính theo glucose). Có tối đa bao nhiêu chiếc gương soi được sản xuất ra? \textit{(Kết quả làm tròn đến hàng đơn vị)}.
\loigiai{}
\end{ex}

\vspace{0.4cm}
\noindent
\textbf{B. PHẦN TỰ LUẬN (3,0 ĐIỂM)}

\setcounter{ex}{0}

\begin{ex}
\textbf{(1,5 điểm)} Viết công thức cấu tạo, gọi tên thay thế và gốc chức của amin bậc 2, bậc 3 có công thức phân tử $\ce{C4H11N}$?
\loigiai{}
\end{ex}

\begin{ex}
\textbf{(1,5 điểm)}
\begin{enumerate}[a.]
	\item So sánh điểm giống và khác nhau giữa xà phòng và chất giặt rửa tổng hợp?
	\item Trình bày cơ chế giặt rửa của chúng?
\end{enumerate}
\loigiai{}
\end{ex}

\begin{center}
\textbf{--- HẾT ---}
\end{center}

\end{document}
